% Français 9115, batch 29 — Gallica views 561–580.
% Pass: Opus 5.5 (claude-opus-5-5), 2026-09-22 — first pass, unchecked against the views by a human.
%
% What the twenty views hold. Eighteen views carry text; two are blank
% backs (v. 562, 576). All of it belongs to the circular plate of § 7 of
% her memoir on vibrating elastic surfaces — the paginated run that
% batches 20–28 followed from page (1) at v. 398 to (139) at v. 560 — with
% three slips bound among its pages. No number theory, no microfilm
% target, no printed matter, no repeated exposure.
%
%   v. 561      a small calculation slip on a mount, ink 276, no page
%               number: from (1/r)dz/dr + d²z/dr² = (u²/e²)z, derived
%               three times, the fourth-order equation of the round
%               plate when z does not depend on the angle, « comprise dans
%               celle qui a été donnée 139 »; then the same derivation
%               begun with φ for the left-hand side. Back blank (v. 562).
%   v. 563–564  a slip headed « addition a la p. 140 », ink 277 (v. 563),
%               written on both sides: θ = 0, all nodal lines annular;
%               the conditions at the edge satisfied by z = 0 or dz/dr = 0,
%               so that the lines of rest satisfy them without being
%               lines of support « dans le sens \uncertain{analytique}
%               attribué a ce mot par Euler »; the real edge requires
%               dz/dr = 0 at ρ = R. On the back (v. 564) the condition in
%               series with g = u²/e², the values of A and A′ that satisfy
%               it, and the equation whose roots r < R give the annular
%               nodal lines, left unfinished.
%   v. 565–574  pages (140)–(149) of the memoir, fair hand with
%               corrections, both sides written. (140) the expansion of
%               d⁴u/dr⁴ with u = r^β s, the equation (B″), s as a power
%               series; with n = 1+2β the general z as r^β times two
%               series in A and A′. (141) ∇²z; the four couples of edge
%               equations « donnés N.o 9 » in z′, z″. (142) the eight
%               equations; the first couple (supported points): sin(βφ+B)
%               = 0 gives 2β points of rest, the ends of β diametral
%               lines of support. (143) the periodicity in φ, like the
%               vibrating string; β = 0 forbids B = 0; the double
%               supposition. (144) the surface supported all round is the
%               stretched drumhead Euler gave « dans le tome X de
%               pétersbourg » (« De motu vibratorio tympanorum », her N.o
%               13), the sounds of the one as the squares of those of the
%               other; N.o 22, the second couple (free contour). (145)–(146)
%               A and A′ as constants; the annular nodal lines are mere
%               lines of rest; diametral and circular lines together or
%               apart on the free plate. (147) the third couple (fixed
%               points) and the fourth, cos(βφ+B) = 0: each sector halved
%               by a line of analytic limit. (148) N.o 23: the values of u
%               cannot be found, but Chladni's figures are assigned to β:
%               figs. 104 and 109 to β = 0 … fig. 102 a to β = 5; with r = a
%               the equation of the elastic ring. (149) Euler's equations
%               for the ring at the end of his memoir on the elastic
%               lamina differ from hers only in writing the arcs as
%               ordinates; the sounds as the squares of 2, 4, 6 …, where
%               Euler gives 4, 9, 16, 25 …; he did not notice that the ring
%               can have only an even number of nodes. The page stops
%               mid-sentence; the sequel is (150) at v. 577.
%   v. 575      a slip of another paper, ink 283, not paginated, in a
%               faster, heavier hand, first person: her health has kept
%               her from studying « le mémoire que Monsieur fourier a eu
%               la complaisance de me prêter »; « Dans le N.o d'aout »
%               M. Poisson answers M. Navier (p. \uncertain{438}) that
%               « l'hypothese » is not admissible, knowing « cette
%               hypothese que j'ai donné en 1811 » only through Navier's
%               citation; his § VIII, p. 189, on the circular plate, and
%               two equations quoted from his pp. 190 and 191. Back blank
%               (v. 576). A note bound between (149) and (150); it
%               interrupts nothing, the sentence of (149) running on to
%               (150).
%   v. 577–580  pages (150)–(153): the ring as a straight lamina with
%               supported ends (an even number of nodes on the ring, an
%               odd one on the lamina); u/e = β when there are only
%               diametral lines, the sounds as the squares of β. N.o 24,
%               Chladni's objection (« p. 118 du traité d'acoustique »)
%               that Euler's theory of the ring is not confirmed: the
%               sounds go as the squares of 3, 5, 7 …, the theory as those
%               of 4, 6, 8 …. Her explanation: the bending of an iron
%               band (Chladni's section VI, p. 113, figs. 37–42, the
%               tuning fork, and his N.o 89) lowers its elasticity, which
%               lowers the sound and draws the nodes together where it is
%               most bent; « la courbure considerée en elle même » cannot
%               alter the coefficient, as wind instruments show. The page
%               (153) stops mid-sentence at the foot of v. 580.
%
% Seams. At 560/561: view 560 (batch 28), page (139), ends on the
% expansion of d³u/dr³; the slips of v. 561 and v. 563–564 are bound
% between (139) and (140), and page (140) at v. 565 picks the expansion up
% at d⁴u/dr⁴ — the text of the memoir runs from 560 to 565 over the
% slips. View 563 calls itself « addition a la p. 140 ». At 580/581: view
% 580 ends « … comme dans la bande » and view 581 (outside the batch,
% glanced at) is page (154), ink 286, opening « droite, et lorsque sa
% courbure est uniforme … »: the sentence runs on.
%
% The numbers on the leaves, and why no \folio{} is written. (1) Her
% pagination, in parentheses at the head, centre: (140) v. 565, (141)
% v. 566 — the middle figure blotted on both, read 4 from the sequence —
% (142) v. 567, (143) v. 568, (144) v. 569, (145) v. 570, (146) v. 571,
% (147) v. 572, (148) v. 573, (149) v. 574, (150) v. 577, (151) v. 578,
% (152) v. 579, (153) v. 580, without a gap. Struck through: (142),
% (144), (146), (148), (150), (152), the even pages, which stand on the
% side that carries the ink number; (140), on such a side too, and the odd
% pages are not struck. (2) The ink series at the top
% right, in the larger rounder hand, on the recto only: 276 (v. 561),
% 277 (563), 278 (565), 279 (567), 280 (569), 281 (571), 282 (573), 283
% (575, the slip), 284 (577), 285 (579); 286 follows at v. 581. One
% number per physical leaf, slips included, none on the backs. Penned,
% as throughout the volume: no \folio{}. (3) Her paragraph numbers: N.o 22
% (v. 569), N.o 23 (v. 573), N.o 24 (v. 577); references to N.o 2
% (v. 566, 573), N.o 9 (v. 566), N.o 13 (v. 569), N.o 22 (v. 577), and to
% her « p. 140 » (v. 563) and « 139 » (v. 561). (4) Chladni's figures
% 99–111, 116 a (v. 573), 37, 38–42 (v. 578); his traité d'acoustique
% pp. 113, 118 and N.o 89; Poisson's pp. 189–191 and § VIII (v. 575).
% Every number was read on its view; none is inferred. None of these
% leaves can be Laubenbacher and Pengelley's « Manuscript C » (ff.
% 348r–349r): the subject is elasticity, and the ink series here runs
% 276–285.
%
% Notation. Hers, kept. u is the frequency parameter, e the constant of
% the equation (bτ² = e⁴/a⁴ at v. 560), n = 1+2β; A, A′ her series
% constants, written as looped cursive capitals; the separate cursive
% constant of v. 573 is set $\mathcal{A}$. Her fractions of fractions
% (u²/e² over 2(n+1)) are kept as she writes them. « sin », « cos »,
% with or without a stop, are \text{sin}, \text{cos}. Her braces are kept,
% but she often opens with « { » and closes with « ) » or with a brace
% and no partner; the long equations are regularised to balanced
% \left\{ \right\} and split over two rows, which the page does not do.
% « &c » is \&c. Words she underlines in the prose are \emph{};
% underlined letters in running text (u, r, z, x, y, β) keep
% \underline{} inside maths. Plurals she cancels with a small cross are
% given in the singular, with a note. The long s is transcribed s.
% Spelling is hers and is not flagged: « coëfficiens », « differens »,
% « proportionels » and « proportionnels », « quarrés », « constans »,
% « instrumens », « fesant », « appuyés » for the ends, « nauds » given
% as nœuds, « resulteroit », « connoître », « paroit », « ceqi », and the
% run-together words « dela », « laplaque », « lamême », « àlafois »,
% « cequi », « endeux », « delignes », « aumonde », « c'est adire »,
% « desorte ». Where the leaf and the calculation disagree the leaf
% stands, with a note (v. 561, 564, 565, 569, 571).
%
% What could not be read, and what is offered with doubt. \ill{} stands
% four times: two exponents in the struck, overwritten first term of the
% last line of v. 564; a struck word after « on trouvera » at v. 565;
% the blotted end of « compar- » at v. 578. \uncertain{} stands for:
% u⁴/e⁴ (v. 561, in a note); « analytique » (v. 563); two r's and a
% denominator (v. 564); a denominator 2, the page figure 4 and the struck
% « ego » (v. 565); the page figure 4 (v. 566); « uu » and « aux »
% (v. 568); a denominator 4 where (n+3) is expected (v. 569); a crossed
% equals sign read ≠ and a blotted letter read α (v. 570); « o » in « N.o
% 2 » (v. 573); « aux » (v. 574); « 438 » and « traiteroit » (v. 575);
% « negl » (struck), « tems », and « ee » in the place of a figure number
% (v. 578).
%
% Nothing here was checked against any published edition of Euler, of
% Chladni, of Poisson or of Sophie Germain's papers. That the « N.o
% d'aout » is of the Annales de chimie et de physique and that the
% exchange is Poisson and Navier's of 1828 is the pass's own inference,
% stated as such in the note to v. 575.
\documentclass[11pt,a4paper]{article}
% The preamble shared by every transcription of the mathematicians' manuscripts in Gallica.
%
% It rests on one idea: **uncertainty compounds, it does not resolve.** A
% transcription that smooths over a doubtful word has destroyed the only thing
% it was made for — a reader must be able to tell, at every point, what was on
% the page and what was supplied. The apparatus macros below are therefore the
% critical apparatus, not decoration.
%
% The same commands are recognised by `scripts/render.mjs`, which produces the
% site's reading view, and by `scripts/tei.mjs`. A macro added here must be
% added there too, or rendering fails loudly — which is the wanted behaviour.
%
% Comments here are English, like the rest of the repository. The documents
% this preamble sets are French, and that is deliberate: see the skills.

\ProvidesPackage{germain}[2026/09/15 Transcriptions of mathematicians' manuscripts in Gallica]

\usepackage{amsmath,amssymb,amsthm}
\usepackage{mathrsfs}
% Kept from the parent preamble so the renderer's diagram support stays
% available; these manuscripts draw no commutative diagrams, and nothing here uses it.
\usepackage{tikz-cd}
\usepackage[english,french]{babel}
\usepackage{fontspec}

% --- Greek ------------------------------------------------------------------
% The text face stays fontspec's default, Latin Modern, which has no Greek:
% XeTeX drops every glyph it lacks with a line in the log and nothing on the
% page, and the Greek words the manuscripts quote vanished from the PDFs. So
% the Greek and Greek Extended blocks get a XeTeX character class of their own,
% and entering or leaving a run of it switches the family to CMU Serif — the
% same Computer Modern design, with polytonic Greek — and back. Only the
% family changes: a Greek word in \textit or \textbf keeps its shape and series.
% The transitions to and from every other class carry the tokens that class
% has towards an ordinary letter, so French spacing before « ; » or inside
% « » still holds after a Greek word. No grouping, so no brace or \endgroup
% can come out unbalanced; maths is left alone, since XeTeX inserts nothing
% there. Kept identical in `exercices.sty`.
\makeatletter
\newfontfamily\germainGreekfont{cmun}[
  Extension=.otf, NFSSFamily=germaingreek,
  UprightFont=*rm, ItalicFont=*ti, SlantedFont=*sl,
  BoldFont=*bx, BoldItalicFont=*bi, BoldSlantedFont=*bl]
\newXeTeXintercharclass\germain@greek
\def\germain@greekrange#1#2{%
  \@tempcnta=#1\relax
  \loop
    \XeTeXcharclass\@tempcnta=\germain@greek
    \ifnum\@tempcnta<#2\advance\@tempcnta\@ne\repeat}
\germain@greekrange{"0370}{"03FF}
\germain@greekrange{"1F00}{"1FFF}
\def\germain@greekprev{\rmdefault}
\def\germain@greekin{%
  \edef\germain@greekprev{\f@family}\fontfamily{germaingreek}\selectfont}
\def\germain@greekout{\fontfamily{\germain@greekprev}\selectfont}
\def\germain@greekjoin{%
  \XeTeXinterchartokenstate=\@ne
  \@tempcnta=\z@
  \loop
    \ifnum\@tempcnta=\germain@greek\else
      \XeTeXinterchartoks\germain@greek\@tempcnta=\expandafter{%
        \expandafter\germain@greekout\the\XeTeXinterchartoks\z@\@tempcnta}%
      \XeTeXinterchartoks\@tempcnta\germain@greek=\expandafter{%
        \the\XeTeXinterchartoks\@tempcnta\z@\germain@greekin}%
    \fi
    \ifnum\@tempcnta<4095\advance\@tempcnta\@ne\repeat}
% babel-french sets its own classes, and saves and restores the interchar
% state, each time a language is selected: join again after it does.
\addto\extrasfrench{\germain@greekjoin}
\addto\extrasenglish{\germain@greekjoin}
\AtBeginDocument{\germain@greekjoin}
\makeatother

\usepackage{geometry}
\geometry{a4paper,margin=25mm,marginparwidth=28mm}
\usepackage{marginnote}
\usepackage{xcolor}
\usepackage[normalem]{ulem}
\setlength{\emergencystretch}{3em}
\usepackage{hyperref}
\hypersetup{colorlinks=true,linkcolor=black,urlcolor=black,pdfborder={0 0 0}}

% --- Batch metadata ---------------------------------------------------------
% Taken as they stand by the rendering script, which shows them at the head of
% the reading view. They fix what the file claims to cover. The macro names are
% the parent project's, so that the scripts stay shared; \folder here means the
% volume's slug — `fr-9115` — and \pages counts Gallica views.

\newcommand{\folder}[1]{\gdef\grothFolder{#1}}
\newcommand{\batch}[1]{\gdef\grothBatch{#1}}
\newcommand{\pages}[2]{\gdef\grothFirst{#1}\gdef\grothLast{#2}}
\newcommand{\foldertitle}[1]{\gdef\grothFoldertitle{#1}}

% The holder's own shelfmark, printed as they write it: « Français 9115 ».
\newcommand{\shelfmark}[1]{\gdef\grothShelfmark{#1}}

% The dating, copied verbatim from the holder's catalogue — never inferred
% here. « XIXe siècle » is the BnF's; a pass that dates a leaf from its content
% says so in a \note{}, as its own inference.
\newcommand{\dating}[1]{\gdef\grothDating{#1}}

% The status notice, printed in the title block and repeated as a diagonal
% watermark on every page of the PDF. The manuscripts are in the public
% domain; what this line declares is not a right but a quality — an
% unverified first pass by a machine. The skills write the line; a file
% without it gets no watermark, so leaving it out is visible in the output.
\newcommand{\watermark}[1]{\gdef\grothWatermark{#1}}

\gdef\grothFolder{?}\gdef\grothBatch{}\gdef\grothFirst{?}\gdef\grothLast{?}%
\gdef\grothFoldertitle{}\gdef\grothShelfmark{}\gdef\grothDating{}\gdef\grothWatermark{}

% --- The résumé -------------------------------------------------------------
\newenvironment{resume}
  {\small\setlength{\parskip}{0.45em}}
  {\par\medskip\hrule\medskip}

% --- Keywords ---------------------------------------------------------------
% In English, closing the modernised reading's résumé: the modern vocabulary
% under which to search for what the leaves construct. The manifest extracts
% them to tag the volume — this line is the single source of the tags.
\newcommand{\keywords}[1]{\par\smallskip\noindent
  {\footnotesize\textsc{Keywords} --- \textit{#1}}\par}

% --- The critical apparatus -------------------------------------------------

% The Gallica view number, in the margin. This is the anchor that lets a
% disagreement over a formula be settled by looking at *one* image rather than
% a volume of 750. The site uses it to turn the facsimile as the transcription
% is scrolled. A view is one scanned image: usually one side of a leaf.
\newcommand{\page}[1]{%
  \marginnote{\footnotesize\color{gray!70}\textsf{v.\,#1}}%
  \label{page:#1}\ignorespaces}

% The BnF's pencilled foliation, where the leaf carries it and a pass can read
% it: « 198r ». This is what the literature cites — « ff. 198r–208v » — and
% the one line that turns a citation into a view. Recorded, never inferred:
% a folio a pass cannot read is not written.
\newcommand{\folio}[1]{%
  \marginnote{\footnotesize\color{gray!50}\textsf{f.\,#1}}\ignorespaces}

% The modernised reading's coarser counterpart to \page{}: which views a
% section is drawn from.
\newcommand{\pagerange}[2]{%
  \marginnote{\footnotesize\color{gray!70}\textsf{v.\,#1--#2}}%
  \ignorespaces}

% Illegible: never guessed.
\newcommand{\ill}{\textcolor{red!60!black}{[\,\ldots\,]}}

% A doubtful reading: the word is given, and flagged as uncertain.
\newcommand{\uncertain}[1]{\uline{#1}}

% An editorial addition — a missing letter, an implied word.
\newcommand{\add}[1]{\textcolor{blue!50!black}{[#1]}}

% Struck out by the writer. What was crossed out is part of the document.
\newcommand{\struck}[1]{\sout{#1}}

% The transcriber's note, on the state of the page or the mathematics.
\newcommand{\note}[1]{\par\smallskip{\small\color{gray!60!black}\textit{[#1]}}\par\smallskip}

% A marginal note of hers, distinct from the body of the text.
\newcommand{\marginal}[1]{\par\smallskip\noindent\hspace{1em}%
  {\small\color{gray!60!black}#1}\par\smallskip}

% --- Title ------------------------------------------------------------------
% The title block is French, like the documents themselves.

\AddToHook{shipout/background}{%
  \ifx\grothWatermark\empty\else
    \begin{tikzpicture}[remember picture, overlay]
      \node[rotate=52, scale=3.4, align=center, text=gray!18,
            font=\sffamily\bfseries]
        at (current page.center) {\grothWatermark};
    \end{tikzpicture}%
  \fi}

\AtBeginDocument{%
  \begin{center}
    {\large\bfseries Manuscrits de mathématiciens (Gallica) ---
      \ifx\grothShelfmark\empty\grothFolder\else\grothShelfmark\fi}\\[2pt]
    {\itshape\grothFoldertitle}\\[4pt]
    {\small vues \grothFirst--\grothLast
      \ifx\grothBatch\empty\else\ (lot \grothBatch)\fi}\\[2pt]
    \ifx\grothDating\empty\else
      {\small\color{gray!60!black}Datation du catalogue : \grothDating}\\[2pt]
    \fi
    \ifx\grothWatermark\empty\else
      {\small\bfseries\color{red!45!gray}\grothWatermark}\\[2pt]
    \fi
    {\footnotesize\color{gray!60!black}Transcription établie d'après les images IIIF de Gallica.
     Source gallica.bnf.fr / Bibliothèque nationale de France.\\
     Les lectures incertaines sont soulignées, les illisibles notées [\,\ldots\,],
     les ajouts entre crochets ; \textsf{v.} numérote les vues de Gallica,
     \textsf{f.} la foliotation au crayon de la BnF.}
  \end{center}
  \vspace{1em}}

\endinput


\folder{fr-9115}
\shelfmark{Français 9115}
\batch{29}
\pages{561}{580}
\dating{1801-1900}
\watermark{Lecture automatique\\non vérifiée}
\foldertitle{Papiers de Sophie GERMAIN. — Recueil de dissertations et problèmes mathématiques et physiques.}

\begin{document}

\note{Numérotation des feuillets. Les pages du mémoire portent en tête,
au milieu, entre parenthèses, leur pagination : (140) à la vue 565, puis
(141) à (149) aux vues 566 à 574, et (150) à (153) aux vues 577 à 580,
sans lacune ; les pages paires (142) à (152) sont barrées, (140) ne
l'est pas. Le côté de chaque feuillet qui porte une page paire porte
aussi, en haut à droite,
à l'encre, d'une écriture plus ronde, un nombre de la série qui court
dans tout le volume : 276 (vue 561), 277 (563), 278 (565), 279 (567),
280 (569), 281 (571), 282 (573), 283 (575), 284 (577), 285 (579) ; les
petits feuillets des vues 561, 563 et 575, sans pagination, en ont un
aussi. Cette série occupe la place de la foliotation de la
Bibliothèque, mais elle est tracée à la plume et non au crayon, et aucun
« f. » n'est donc porté. Tous ces nombres ont été lus sur la vue ; aucun
n'est déduit. Les vues 562 et 576 sont les dos blancs des feuillets des
vues 561 et 575 et ne sont pas transcrites.}

\page{561}

On a
\[ \frac{1}{r}\frac{dz}{dr} + \frac{d^2z}{dr^2} = \frac{u^2}{e^2}z \]
\[ -\frac{1}{r^2}\frac{dz}{dr} + \frac{1}{r}\frac{d^2z}{dr^2} + \frac{d^3z}{dr^3} = \frac{u^2}{e^2}\frac{dz}{dr} \]
\[ \frac{2}{r^3}.\frac{dz}{dr} - \frac{2\,d^2z}{r^2dr^2} + \frac{1}{r}\frac{d^3z}{dr^3} + \frac{d^4z}{dr^4} = \frac{u^2}{e^2}\frac{d^2z}{dr^2} \]
\[ \frac{2}{r^3}.\frac{dz}{dr} - \frac{2\,d^2z}{r^2dr^2} + \frac{1}{r}\frac{d^3z}{dr^3} + \frac{d^4z}{dr^4} + \frac{1}{r^3}\frac{dz}{dr} + \frac{1}{r^2}\frac{d^2z}{dr^2} + \frac{1}{r}\frac{d^3z}{dr^3} = \frac{u^2}{e^2}\left(\frac{1}{r}\frac{dz}{dr} + \frac{d^2z}{dr^2}\right) \]
En reduisant on trouve
\[ \frac{1}{r^3}\frac{dz}{dr} - \frac{d^2z}{r^2dr^2} + \frac{2\,d^3z}{r\,dr^3} + \frac{d^4z}{dr^4} = \frac{u^4}{e^4}z \]
qui est en effet l'équation dela surface
comprise dans celle qui a été donnée 139 et
qui appartient au cas ou l'ordonnée $z$ n'est
pas fonction dela quantité angulaire $\theta$.

\[ \frac{1}{r}\left(\frac{dz}{dr}\right) + \frac{ddz}{dr^2} = \varphi \]
\[ -\frac{1}{r^2}.\frac{dz}{dr} + \frac{1}{r}\frac{d^2z}{dr^2} + \frac{d^3z}{dr^3} = \frac{d\varphi}{dr} \]
\[ +\frac{2}{r^3}\frac{dz}{dr} - \frac{1}{r^2}\frac{d^2z}{dr^2} - \frac{1}{r^2}\frac{d^2z}{dr^2} + \frac{1}{r}\frac{d^3z}{dr^3} + \frac{d^4z}{dr^4} = \frac{d^2\varphi}{dr^2} \]
\note{Petit feuillet de calculs, monté sur un grand feuillet blanc ; en
haut à droite, à l'encre, « 276 ». Aucune pagination de sa main. Dans
la première équation, le dénominateur de $\frac{u^2}{e^2}$ est empâté
d'une tache ; l'exposant du second membre de l'équation réduite
(\uncertain{$u^4/e^4$}) l'est aussi. Au quatrième membre de la
quatrième équation, qui ajoute à la troisième la deuxième divisée par
$r$, le feuillet porte « $+\frac{1}{r^3}\frac{dz}{dr}$ » où le calcul
demande $-\frac{1}{r^3}\frac{dz}{dr}$ ; la réduction qui suit,
$\frac{1}{r^3}\frac{dz}{dr}$, est faite avec le signe moins ; on laisse
le signe du feuillet. Dans le dernier terme de la même ligne, le
dénominateur « $r$ » est suivi d'une tache. « 139 » renvoie
vraisemblablement à la page (139) de sa pagination (vue 560, hors de
ce lot), où l'équation (B) est donnée : c'est l'inférence du
transcripteur. La lettre angulaire finale est une
boucle proche d'un « 8 », lue $\theta$. Sous un trait, un second
calcul : dans la dernière ligne, le premier signe moins est un trait
épais et empâté, et les exposants de $d^3z$ et de $d^4z$ sont récrits
sur d'autres chiffres. Le bas du feuillet est blanc ; le dos (vue 562)
ne montre que l'écriture du recto par transparence.}

\page{563}

addition a la p. 140

Si on prend $\theta = 0$ on aura les cas devibrations
ou les lignes nodales sont toutes annulaires.

Les conditions des extremités sont alors
\[ \left(\frac{1}{r}\left(\frac{dz}{dr}\right) + \frac{ddz}{dr^2}\right)\delta\frac{dz}{dr}\,\text{cos}.\varphi + \left(\frac{1}{r^2}\frac{dz}{dr} - \frac{1}{r}\frac{ddz}{dr^2} - \frac{d^3z}{dr^3}\right)\delta z\,\text{cos}.\varphi = 0 \]
et
\[ \left(\frac{1}{r}\left(\frac{dz}{dr}\right) + \frac{ddz}{dr^2}\right)\delta\frac{dz}{dr}\,\text{sin}.\varphi + \left(\frac{1}{r^2}\frac{dz}{dr} - \frac{1}{r}\frac{ddz}{dr^2} - \frac{d^3z}{dr^3}\right)\delta z\,\text{sin}.\varphi = 0 \]
A cause de
\[ \frac{1}{r}\left(\frac{dz}{dr}\right) + \frac{ddz}{dr^2} = \frac{e^2}{a^2}z \]
\[ -\frac{1}{r^2}\left(\frac{dz}{dr}\right) + \frac{1}{r}\frac{d^2z}{dr^2} + \frac{d^3z}{dr^3} = \frac{e^2}{a^2}\frac{dz}{dr} \]
on peut y satisfaire par l'une ou l'autre des suppositions
\[ z = 0 \qquad \frac{dz}{dr} = 0. \]
On voit parla que toutes les lignes de repos satisfont
aux conditions des extremités sans être pourtant des
lignes d'appuis dans le sens \uncertain{analytique} attribué a ce
mot par Euler: on voit encore que toutes les fois
que la tangente est parallele auplan dela surface
en repos on a une ligne circulaire qui satisfait
aussi aux conditions des extremités. Ces deux remarques
établissent entre le cas que nous examinons et
ceux qui appartiennent aux intégrales angulaires
et exponentielles, c'est adire tous ceux qui ont été
examinés jusqu'ici, une difference remarquable.
Les extremités réelles et physiques sont assujetties a
la condition $\frac{dz}{dr} = 0$ ainsi elle a lieu lorsque
le rayon $\rho$ a obtenu savaleur complette $R$.
\note{Petit feuillet monté sur un grand feuillet blanc, écrit sur le
seul quart supérieur de la vue ; en haut à droite, à l'encre, « 277 ».
Aucune pagination de sa main : la ligne de tête, « addition a la
p. 140 », dit que la feuille est un ajout à la page (140) de sa
pagination, qui est la vue 565. Le « 4 » est tracé comme un « b », à
la manière qu'on lui connaît, et pourrait passer pour un 6 ; la page
(140) de la vue 565, qui reprend le calcul là où la page (139) de la
vue 560 l'a laissé, décide. La lettre angulaire de la
première ligne est la même boucle qu'à la vue 561, lue $\theta$ ; la
parenthèse qui ferme le premier facteur des deux équations des
extrémités est tracée comme une accolade. Dans la seconde équation
« à cause de », l'exposant de $a$ est mal formé. Le mot lu
\uncertain{analytique} s'écrit sur la vue « ahotytique » ou
approchant : la seconde lettre est une haste bouclée, suivie d'un
« o » ; le sens appelle « analytique », qu'on n'impose pas. « Euler »
est écrit avec l'E en forme de L. Le rayon est d'abord $r$ dans les
équations, puis $\rho$ dans la dernière ligne ; on garde les deux. Un
timbre rond de la Bibliothèque (« Bibliothèque impériale MSS. ») est
imprimé à droite des équations ; deux gouttes d'encre sont tombées
sur la ligne « dans le sens … ». Le dos (vue 564) porte des calculs.}

\page{564}

En mettant $g$ alaplace de $\frac{u^2}{e^2}$ cette condition donne l'équation
\[ A'\left\{\frac{gR}{2} + \frac{g^2R^3}{2.2.4} + \frac{g^3R^5}{2.4.2.4.6} + \frac{g^4R^7}{2.4.6.2.4.6.8} + \&c\right\} \]
\[ -A\left\{\frac{gR}{2} - \frac{g^2R^3}{2.2.4} + \frac{g^3R^5}{2.4.2.4.6} - \frac{g^4R^7}{2.4.6.2.4.6.8} + \&c\right\} = 0 \]
$R$ est une quantité constante parconséquent on satisfera
a cette equation en prenant
\[ A = 1 + \frac{gR^2}{2.4} + \frac{g^2R^4}{4.2.4.6} + \frac{g^3R^6}{4.6.2.4.6.8} + \&c. \]
\[ A' = 1 - \frac{gR^2}{2.4} + \frac{g^2R^4}{4.2.4.6} - \frac{g^3R^6}{4.6.2.4.6.8} + \&c. \]
En général il y aura autant deligne nodales annulaires
sur la plaque ronde qu'il y a de valeur de $r < R$
qui pour une valeur donnée de $g$ satisferont a la condition
\[ \left(1 + \frac{gR^2}{2.4} + \frac{g^2R^4}{4.2.4.6} + \frac{g^3R^6}{4.6.2.4.6.8} + \&c.\right) \]
\[ \left\{1 - \frac{g\uncertain{r}^2}{(2)^2} + \frac{g^2\uncertain{r}^4}{(2.4)^2} - \frac{g^3r^6}{(2.4.6)^2} + \&c\right) \]
\[ +\left(1 - \frac{gR^2}{2.4} + \frac{g^2R^4}{4.2.4.6} - \frac{g^3R^6}{4.6.2.4.6.8} + \&c\right) \]
\[ \left\{1 + \frac{gr^2}{(2)^2} + \frac{g^2r^4}{(2.4)^2} + \frac{g^3r^6}{(2.4.6)^2} + \&c\right) = 0 \]
En effaçant cequi se détruit cette equation \struck{peut} devient
\[ 1 + \frac{g^2R^4}{4.2.4.6} + \frac{g^4R^8}{4.6.8.2.4.6.8.10} + \&c. \]
\[ -\frac{gR^2}{2.4}.\frac{gr^2}{(2)^2} - \frac{g^3R^6}{4.6.2.4.6.8}.\frac{gr^2}{(2)^2} - \&c. \]
\[ +\frac{g^2r^4}{(2.4)^2} + \frac{g^2R^4}{4.2.4.6}.\frac{g^2r^4}{(2.4)^2} + \frac{g^4R^8}{4.6.8.2.4.6.8.10}\frac{g^2r^4}{(2.4)^2} + \&c.\Big) = \]
\[ -\frac{g^{\ill{}}R^{\ill{}}}{\uncertain{4.6.2.4.6.8}}.\frac{g^3r^6}{(2.4.6)^2} - \frac{g^3R^6}{4.6.2.4.6.8}\frac{g^3r^6}{(2.4.6)^2} - \&c \]
\[ + \&c \]
\note{Dos du feuillet de la vue 563 (l'écriture de l'un se voit à
l'envers par transparence sur l'autre) ; pas de nombre en haut à
droite. Les accolades et parenthèses sont les siennes, et elle ouvre
souvent par « $\{$ » ce qu'elle ferme par « ) ». Dans le premier facteur
de la seconde ligne de la condition, les deux premiers $r$ sont
couverts d'une tache et lus par la suite du calcul ; le signe entre
les deux derniers termes du facteur suivant est un trait épais. Une
grande parenthèse, en marge droite, embrasse les trois lignes qui
suivent « devient », et un signe « = » est écrit tout au bord du
feuillet après elle ; rien ne suit sur la vue. Dans la dernière ligne,
le premier terme est surchargé : les exposants du numérateur sont
empâtés (la lettre lue $R$), le dénominateur a le premier chiffre et
« 6.8 » barrés ou récrits ; au facteur suivant, la lettre du
numérateur est tracée comme un « p », pour $r$. Le « + » devant le dernier « \&c. » de l'avant-dernière ligne
est traversé d'un trait.}

\page{565}

\[ \frac{d^4u}{dr^4} = \beta(\beta-1)(\beta-2)(\beta-3)\,r^{\beta-4}s + 4\beta(\beta-1)(\beta-2)\,r^{\beta-3}.\frac{ds}{dr} + 6\beta(\beta-1)\,r^{\beta-2}.\frac{dds}{dr^2} + 4\beta\,r^{\beta-1}.\frac{d^3s}{dr^3} \]
\[ + \frac{d^4s}{dr^4}. \]
\struck{et} apres avoir substitué ces differentes valeurs dans l'équation
$(B')$, \struck{et} \struck{avoir} \add{avoir} pratiqué les reductions convenables \struck{on trouvera}
\struck{aura enfin l'équation} et divisé par $r^\beta$ on aura enfin
l'équation
\[ \frac{u^4}{e^4}s = \frac{d^4s}{dr^4} + (1+2\beta)\left\{\frac{2}{r}.\frac{d^3s}{dr^3} + (2\beta-1)\left\{\frac{1}{r^2}.\frac{dds}{dr^2} - \frac{1}{r^3}.\frac{ds}{dr}\right\}\right\} \ldots (B'') \]
La valeur la plus générale que l'on puisse attribuer à $s$ est
$s = A + Br + Cr^2 + Dr^3 + \&c.$; car si on \struck{vouloit} prenoit
\[ s = \text{sin}.\frac{u}{e}r\left\{a + br + cr^2 + dr^3 + er^4 + \&c\right\} + \text{cos}.\frac{u}{e}r\left\{A + Br + Cr^2 + Dr^3 + Er^4 + \&c\right\} \]
Le developpement des valeurs de $\text{sin}.\frac{u}{e}r$ et $\text{cos}.\frac{u}{e}r$ en serie
rameneroit à la première \struck{valeur} \add{forme attribuée à} \struck{pour} $\underline{s}$.

Pour déterminer les coëfficiens des differentes puissances de $\underline{r}$
dans le developpement dela valeur de $\underline{s}$, il faudra,
\struck{\uncertain{ego}} après avoir substitué \struck{dans l'équation $(B'')$} aux différentielles
de \struck{$d$} $\underline{s}$ leurs valeurs en $\underline{r}$, égaler à zéro la somme
de chacun des coëfficiens \struck{de chacune} des \add{différentes} puissances de
$\underline{r}$ dans l'équation $(B'')$ \struck{ainsi transformée en suivant et}
on trouvera \struck{\ill{}} ainsi, en fesant $1+2\beta = n$
\[ s = A\left\{1 - \frac{\frac{u^2}{e^2}r^2}{\uncertain{2}(n+1)} + \frac{\frac{u^4}{e^4}.r^4}{2.4(n+1)(n+3)} - \frac{\frac{u^6}{e^6}r^6}{2.4.6.(n+1)(n+3)(n+5)} + \&c\right\} \]
\[ + A'\left\{1 + \frac{\frac{u^2}{e^2}r^2}{2(n+1)} + \frac{\frac{u^4}{e^4}.r^4}{2.4(n+1)(n+3)} + \frac{\frac{u^6}{e^6}r^6}{2.4.6(n+1)(n+3)(n+5)} + \&c\right\} \]
formule dans laquelle $A$ et $A'$ sont deux constantes
arbitraires.
\note{En tête, au milieu, entre parenthèses, « (1\uncertain{4}0) » : le
chiffre du milieu est couvert d'une tache ; il est lu 4 parce que la
page (139), à la vue 560 (hors de ce lot), s'arrête sur
$\frac{d^3u}{dr^3}$ et que celle-ci commence par $\frac{d^4u}{dr^4}$,
et parce que le feuillet de la vue 563 se dit « addition a la p. 140 ».
En haut à droite, à l'encre, « 278 ». Dans le développement de
$\frac{d^4u}{dr^4}$, le dernier terme est écrit $\frac{d^4s}{dr^4}$ sans
le facteur $r^\beta$ que portent les autres ; on laisse le feuillet.
Dans l'équation (B$''$), après « $2\beta$ » une lettre est récrite et
empâtée ; les deux accents de « B$''$ » ressemblent à un petit « 4 ».
Le premier mot barré avant « aura enfin l'équation » pourrait se lire
« l'aura ». « prenoit » est traversé par le prolongement du trait qui
barre « vouloit » ; on le garde, la phrase le demande. « forme
attribuée à » est écrit au-dessus de « valeur », barré. Le mot barré
devant « ainsi », après « on trouvera », ne se lit pas ; le premier
mot de la ligne « après avoir substitué » est barré et lu, avec doute,
« ego ». Les capitales $A$, $A'$, $B$, $C$, $D$, $E$ sont des capitales
cursives bouclées. Le premier dénominateur, lu $2(n+1)$ comme le
second, a un chiffre bouclé qui peut passer pour un 8. Le 5 de
$(n+5)$ a une queue descendante. Les lettres soulignées sont
soulignées sur la page. Le verso de ce feuillet est la vue 566.}

\page{566}

et la valeur de $\underline{z}$ sera \struck{donnée par l'équation} donc
\[ z = \ldots\,\text{sin}(\beta\varphi + B)r^\beta\left\{\begin{array}{l} A\left\{1 - \frac{\frac{u^2}{e^2}.r^2}{2(n+1)} + \frac{\frac{u^4}{e^4}.r^4}{2.4(n+1)(n+3)} - \frac{\frac{u^6}{e^6}.r^6}{2.4.6(n+1)(n+3)(n+5)} + \&c\right\} \\ + A'\left\{1 + \frac{\frac{u^2}{e^2}r^2}{2(n+1)} + \frac{\frac{u^4}{e^4}.r^4}{2.4(n+1)(n+3)} + \frac{\frac{u^6}{e^6}.r^6}{2.4.6(n+1)(n+3)(n+5)} + \&c\right\} \end{array}\right. \]
et on \struck{trouvera} aura, après les réductions convenables,
\[ \frac{ddz}{dx^2} + \frac{ddz}{dy^2} = \frac{u^2}{e^2}r^\beta\,\text{sin}.(\beta\varphi + B)\left\{\begin{array}{l} A'\left\{1 + \frac{\frac{u^2}{e^2}r^2}{2(n+1)} + \frac{\frac{u^4}{e^4}r^4}{2.4.(n+1)(n+3)} + \&c\right\} \\ - A\left\{1 - \frac{\frac{u^2}{e^2}r^2}{2(n+1)} + \frac{\frac{u^4}{e^4}.r^4}{2.4.(n+1)(n+3)} - \&c\right\} \end{array}\right. \]
\struck{En général} Si on représente par $z'$ et $z''$ les valeurs de $\underline{z}$
qui satisfont \struck{au} respectivement aux deux équations du
second degré trouvées N.o 2, \struck{il est que} ces valeurs seront
ici
\[ z' = r^\beta A\left\{1 - \frac{\frac{u^2}{e^2}r^2}{2(n+1)} + \frac{\frac{u^4}{e^4}r^4}{2.4(n+1)(n+3)} - \&c\right\} \]
\[ z'' = r^\beta A'\left\{1 + \frac{\frac{u^2}{e^2}r^2}{2(n+1)} + \frac{\frac{u^4}{e^4}r^4}{2.4(n+1)(n+3)} + \&c\right\} \]
et les quatre couples d'équations donnés N.o 9 comme devant
appartenir aux points du contour savoir:
\[ z'-z''=0 \text{ et } z+z''=0\,; \quad z-z''=0 \text{ et } \frac{dz'}{d\varphi} - \frac{dz''}{d\varphi} = 0\,; \]
\[ z'+z''=0 \text{ et } \frac{dz'}{d\varphi} + \frac{dz''}{d\varphi} = 0\,; \quad \frac{dz'}{d\varphi} - \frac{dz''}{d\varphi} = 0 \text{ et } \frac{dz'}{d\varphi} + \frac{dz''}{d\varphi} = 0. \]
Deviendront après avoir été divisées par $r^\beta$,
qui ne peut jamais être nul.
\note{En tête, au milieu, « (1\uncertain{4}1) », le chiffre du milieu
empâté comme à la page précédente. Pas de nombre en haut à droite :
c'est le dos du feuillet de la vue 565. Les accolades qui embrassent
les deux lignes des deux grandes formules sont les siennes, sans
accolade fermante. Le signe devant le second $A$ de la deuxième
formule est une tache, lue « $-$ » d'après les signes de la série
qu'il multiplie. Dans le premier des quatre couples, le deuxième $z$
n'a pas son accent en haut, et le troisième $z$ non plus ; une petite
virgule est sous chacun d'eux : on les laisse sans accent, comme sur
la vue. « divisées » : la finale « ées » est écrite au-dessus d'une
autre, empâtée. La phrase « Deviendront … » n'a pas de suite sur la
page : le bas du feuillet est blanc. Le $\underline{z}$ souligné est
souligné sur la page.}

\page{567}

\[ \begin{array}{l} \left\{A\left\{1 - \frac{\frac{u^2}{e^2}r^2}{2(n+1)} + \frac{\frac{u^4}{e^4}r^4}{2.4(n+1)(n+3)} - \&c\right\}\right. \\ \qquad \left. - A'\left\{1 + \frac{\frac{u^2}{e^2}r^2}{2(n+1)} + \frac{\frac{u^4}{e^4}r^4}{2.4(n+1)(n+3)} + \&c\right\}\right\}\text{sin}(\beta\varphi + B) = 0 \end{array} \]
\[ \begin{array}{l} \text{et }\left\{A\left\{1 - \frac{\frac{u^2}{e^2}r^2}{2(n+1)} + \frac{\frac{u^4}{e^4}r^4}{2.4(n+1)(n+3)} - \&c\right\}\right. \\ \qquad \left. + A'\left\{1 + \frac{\frac{u^2}{e^2}r^2}{2(n+1)} + \frac{\frac{u^4}{e^4}r^4}{2.4(n+1)(n+3)} + \&c\right\}\right\}\text{sin}(\beta\varphi + B) = 0; \end{array} \]
\[ \begin{array}{l} \left\{A\left\{1 - \frac{\frac{u^2}{e^2}r^2}{2(n+1)} + \frac{\frac{u^4}{e^4}r^4}{2.4(n+1)(n+3)} - \&c\right\}\right. \\ \qquad \left. - A'\left\{1 + \frac{\frac{u^2}{e^2}r^2}{2(n+1)} + \frac{\frac{u^4}{e^4}r^4}{2.4(n+1)(n+3)} + \&c\right\}\right\}\text{sin}(\beta\varphi + B) = 0 \end{array} \]
\[ \begin{array}{l} \text{et }\left\{A\left\{1 - \frac{\frac{u^2}{e^2}r^2}{2(n+1)} + \frac{\frac{u^4}{e^4}r^4}{2.4(n+1)(n+3)} - \&c\right\}\right. \\ \qquad \left. - A'\left\{1 + \frac{\frac{u^2}{e^2}r^2}{2(n+1)} + \frac{\frac{u^4}{e^4}r^4}{2.4(n+1)(n+3)} + \&c\right\}\right\}\text{cos}(\beta\varphi + B) = 0; \end{array} \]
\[ \begin{array}{l} \left\{A\left\{1 - \frac{\frac{u^2}{e^2}r^2}{2(n+1)} + \frac{\frac{u^4}{e^4}r^4}{2.4(n+1)(n+3)} - \&c\right\}\right. \\ \qquad \left. + A'\left\{1 + \frac{\frac{u^2}{e^2}r^2}{2(n+1)} + \frac{\frac{u^4}{e^4}r^4}{2.4(n+1)(n+3)} + \&c\right\}\right\}\text{sin}(\beta\varphi + B) = 0 \end{array} \]
\[ \begin{array}{l} \text{et }\left\{A\left\{1 - \frac{\frac{u^2}{e^2}r^2}{2(n+1)} + \frac{\frac{u^4}{e^4}r^4}{2.4(n+1)(n+3)} - \&c\right\}\right. \\ \qquad \left. + A'\left\{1 + \frac{\frac{u^2}{e^2}r^2}{2(n+1)} + \frac{\frac{u^4}{e^4}r^4}{2.4(n+1)(n+3)} + \&c\right\}\right\}\text{cos}(\beta\varphi + B) = 0 \end{array} \]
\[ \begin{array}{l} \left\{A\left\{1 - \frac{\frac{u^2}{e^2}r^2}{2(n+1)} + \frac{\frac{u^4}{e^4}r^4}{2.4(n+1)(n+3)} - \&c\right\}\right. \\ \qquad \left. - A'\left\{1 + \frac{\frac{u^2}{e^2}r^2}{2(n+1)} + \frac{\frac{u^4}{e^4}r^4}{2.4(n+1)(n+3)} + \&c\right\}\right\}\text{cos}(\beta\varphi + B) = 0 \end{array} \]
\[ \begin{array}{l} \text{et }\left\{A\left\{1 - \frac{\frac{u^2}{e^2}r^2}{2(n+1)} + \frac{\frac{u^4}{e^4}r^4}{2.4(n+1)(n+3)} - \&c\right\}\right. \\ \qquad \left. + A'\left\{1 + \frac{\frac{u^2}{e^2}r^2}{2(n+1)} + \frac{\frac{u^4}{e^4}r^4}{2.4(n+1)(n+3)} + \&c\right\}\right\}\text{cos}(\beta\varphi + B) = 0; \end{array} \]
Le premier couple appartient aux points d'extremités appuyés.
On peut y satisfaire de plusieurs manières: et, d'abord, en
fesant $\text{sin}(\beta\varphi + B) = 0$; cette condition montre que lorsque $B = 0$,
c'est-adire lorsque le plan des $x$ et $y$ se confond avec
celui dela plaque en repos, il y a $2\beta$ points de repos à la
circonference. Il est clair que la condition dont il s'agit
est entièrement indépendante de la valeur de $\underline{r}$, et que, parconséquent,
les $2\beta$ points de repos de la circonference ne sont autre chose
que les extremités de $\beta$ lignes \struck{nodales} \add{d'appui} diametrales; \struck{ces lignes}
ces lignes satisferont donc dans toute leur etendue aux conditions
des limites, et elles pourroient également exister sur les surfaces tendues.
\note{En tête, au milieu, « (142) », barré d'un trait horizontal ; en
haut à droite, à l'encre, « 279 ». Le texte est écrit sur la partie
basse d'une feuille plus grande, montée sur un onglet ; le haut de la
vue est blanc. La page commence par les huit équations annoncées au
bas de la page (141) (« Deviendront … ») : la phrase court de la vue
566 à celle-ci. Chaque équation ouvre par une accolade et ferme ses
séries par des accolades doubles ou simples, tracées tantôt « $\{$ »,
tantôt « ) » ; elles sont régularisées ici. Le signe entre les deux
séries de la septième équation est empâté, lu « $-$ » par le couple
auquel elle appartient. Les exposants des numérateurs
$\frac{u^2}{e^2}$, $\frac{u^4}{e^4}$ sont écrits au-dessus du $u$ et
repris au dénominateur. « d'appui » est écrit au-dessus de « nodales »,
barré. « appuyés » : la finale est surchargée. Le $\underline{r}$ est
souligné sur la page.}

\page{568}

Les surfaces circulaires qui sont comprises dans la présente intégrale
jouissent, comme on voit, d'une sorte de périodicité relative
aux valeurs de l'angle $\varphi$; \struck{et en effet} \add{Car} la valeur de $\underline{z}$
donne alors, entre $t$ et $\varphi$, l'équation $\frac{ddz}{dt^2} = \frac{ddz}{d\varphi^2}.\frac{\uncertain{uu}}{\beta\beta}$
dela forme de celle \uncertain{aux} cordes vibrantes, et en effet les
differens secteurs formés par les lignes nodales diamétrales
peuvent être superposés, puisqu'ils sont parfaitement égaux
entr'eux.

Lorsque $\beta = 0$, c'est-a-dire lorsqu'il n'y a pas de lignes
nodales diamétrales, on ne peut faire $B = 0$; \struck{par} car il
resulteroit de cette supposition que laplaque ne pourroit
se mouvoir. Le plan des $\underline{x}$ et $\underline{y}$, ne peut donc se confondre
alors avec celui dela plaque en repos.

On peut encore satisfaire au premier couple d'équations
par la double supposition
\[ \begin{array}{l} A\left\{1 - \frac{\frac{u^2}{e^2}r^2}{2(n+1)} + \frac{\frac{u^4}{e^4}r^4}{2.4(n+1)(n+3)} - \&c\right\} \\ \qquad - A'\left\{1 + \frac{\frac{u^2}{e^2}r^2}{2(n+1)} + \frac{\frac{u^4}{e^4}.r^4}{2.4(n+1)(n+3)} + \&c\right\} = 0 \end{array} \]
\[ \begin{array}{l} A\left\{1 - \frac{\frac{u^2}{e^2}r^2}{2(n+1)} + \frac{\frac{u^4}{e^4}.r^4}{2.4(n+1)(n+3)} - \&c\right\} \\ \qquad + A'\left\{1 + \frac{\frac{u^2}{e^2}.r^2}{2(n+1)} + \frac{\frac{u^4}{e^4}.r^4}{2.4(n+1)(n+3)} + \&c\right\} = 0 \end{array} \]
qui se reduit à
\[ A\left\{1 - \frac{\frac{u^2}{e^2}.r^2}{2(n+1)} + \frac{\frac{u^4}{e^4}r^4}{2.4.(n+1)(n+3)} - \&c\right\} = 0 \]
\[ A'\left\{1 + \frac{\frac{u^2}{e^2}.r^2}{2(n+1)} + \frac{\frac{u^4}{e^4}.r^4}{2.4(n+1)(n+3)} + \&c\right\} = 0 \]
Si on fait $A' = 0$ et $1 - \frac{\frac{u^2}{e^2}.r^2}{2(n+1)} + \frac{\frac{u^4}{e^4}r^4}{2.4.(n+1)(n+3)} - \&c = 0$
on aura le cas où la surface seroit appuyée
\note{En tête, au milieu, « (143) », le 4 tracé comme un « b » ; pas
de nombre en haut à droite : c'est le dos du feuillet de la vue 567.
« Car » est écrit au-dessus de « et en effet », barré. La lettre lue
$t$ dans « entre $t$ et $\varphi$ » et dans $dt^2$ est une petite croix
qui peut passer pour un « + ». Le numérateur du dernier facteur, lu
\uncertain{uu}, est deux petites lettres jumelles au-dessus de
$\beta\beta$. Le mot lu \uncertain{aux} devant « cordes » porte une
croix sur sa fin. Dans « on ne peut faire $B = 0$ », la capitale est
empâtée ; c'est la constante $B$ de $\text{sin}(\beta\varphi + B)$, non
$\beta$. « pourroit » : une lettre est récrite au milieu du mot. La
phrase « on aura le cas où la surface seroit appuyée » continue à la
page suivante. Les lettres soulignées le sont sur la page.}

\page{569}

dans tous les points de son contour; son équation \struck{seroit
alors la mem} conviendroit également alors àla surface
tendue demême forme, et elle seroit en effet la même
que Euler adonnée dans le tome X de pétersbourg. (V. N.o 13
du m. \emph{De motu vibratorio tympanorum}). Les figures
nodales seroient parconséquent semblables, sur les deux
genres de surfaces; et, conformement àl'ordre des
equations differentielles qui leur appartiennent \struck{reciproquement} \add{respectivement},
les sons, qui pour l'une d'elles, seroient proportionels aux
valeurs de $\underline{u}$, \struck{seroient par rapport a l'autre} suivroient
pour l'autre la proportion des quarrés des mêmes
valeurs.

N.o 22 Le second couple d'équations convient \add{dans le cas} où les points de
contour sont libres en vertu des conditions $\frac{ddz}{dx^2} + \frac{ddz}{dy^2} = 0$
et $\frac{d\left\{\frac{ddz}{dx^2} + \frac{ddz}{dy^2}\right\}}{dx} = 0$; ou, cequi est ici lamême chose, $\frac{ddz}{dx^2} + \frac{ddz}{dy^2} = 0$
et $\frac{d\left\{\frac{ddz}{dx^2} + \frac{ddz}{dy^2}\right\}}{dy} = 0$, on peut y satisfaire par la supposition
\[ \begin{array}{l} A\left\{1 - \frac{\frac{u^2}{e^2}r^2}{2(n+1)} + \frac{\frac{u^4}{e^4}.r^4}{2.4(n+1)(n+3)} - \&c\right\} \\ \qquad - A'\left\{1 + \frac{\frac{u^2}{e^2}r^2}{2(n+1)} + \frac{\frac{u^4}{e^4}.r^4}{2.4(n+1)(n+\uncertain{4})} + \&c\right\} = 0 \end{array} \]
Lorsque les valeurs de $A$ et $A'$ seront données, en
\struck{mettant} fesant dans cette équation $r = a$, c'est àdire
$r$ égal au rayon du cercle dont la circonférence est
à la limite de la surface, elle pourra servir à
déterminer les differentes valeurs de $\underline{u}$ qui correspondront
aux mouvements dont la plaque \struck{est} \add{sera} susceptible.
\note{En tête, au milieu, « (144) », barré obliquement ; en haut à
droite, à l'encre, « 280 ». La première ligne achève la phrase de la
vue 568 (« on aura le cas où la surface seroit appuyée / dans tous les
points de son contour »). « Euler » est écrit avec l'E en forme de L
et une fin très cursive (« Lobr »), comme aux vues 419–420 ; « tome
X » : le chiffre romain est une croix. « du m. » abrège sans doute
« du mémoire ». Le titre latin est souligné sur la page, en trois
traits. Dans « respectivement », écrit au-dessus de
« reciproquement » barré, la virgule qui suit est la sienne. Au
second couple, le « $dx^2$ » du premier terme de « ou, cequi est ici
lamême chose » est écrit au-dessus d'un autre dénominateur barré. Au
dernier dénominateur de l'équation, le feuillet porte un chiffre lu
\uncertain{4} là où les autres dénominateurs ont $(n+3)$. Dans
« correspondront » la fin est récrite sur une autre ; dans
« mouvements » le « t » est ajouté. Le bas du feuillet est blanc.}

\page{570}

\struck{Mais} \struck{Quelque} \add{quelles que} soient les valeurs de $\underline{u}$, il est visible que
\[ 1 - \frac{\frac{u^2}{e^2}.a^2}{2(n+1)} + \frac{\frac{u^4}{e^4}.a^4}{2.4.(n+1)(n+3)} - \&c \;\uncertain{\neq}\; \uncertain{\alpha} \quad \text{et} \]
\[ 1 + \frac{\frac{u^2}{e^2}.a^2}{2(n+1)} + \frac{\frac{u^4}{e^4}.a^4}{2.4.(n+1)(n+3)} + \&c \;\uncertain{\neq}\; \uncertain{\alpha} \]
seront des quantités constantes; \struck{ainsi} \add{donc} on pourra toujours
prendre,
\[ A = 1 - \frac{\frac{u^2}{e^2}.a^2}{2(n+1)} + \frac{\frac{u^4}{e^4}.a^4}{2.4(n+1)(n+3)} - \&c \]
\[ A' = 1 + \frac{\frac{u^2}{e^2}.a^2}{2(n+1)} + \frac{\frac{u^4}{e^4}.a^4}{2.4(n+1)(n+3)} + \&c. \]
\add{Mais dans cette supposition} \struck{Et} \struck{alors} les conditions des extremités auront lieu d'elles
mêmes, et ne pourront plus servir, parconséquent à connoître
les valeurs de $\underline{u}$ \struck{puisque} qui devront être donnés d'ailleurs.

Le cas dont je viens de parler, sera celui où dans
la plaque circulaire, on n'auroit pas égard àla
variabilité du rayon; ce cas est absolument analogue
à celui où, dans les plaques rectangulaires, on
ne regarde qu'une des deux coordonnés $\underline{x}$ et $\underline{y}$ comme
variable.

En général, si pour une valeur quelconque de $\underline{u}$,
on peut satisfaire à l'équation
\[ \begin{array}{l} A\left\{1 - \frac{\frac{u^2}{e^2}.r^2}{2(n+1)} + \frac{\frac{u^4}{e^4}.r^4}{2.4(n+1)(n+3)} - \&c\right\} \\ \qquad + A'\left\{1 + \frac{\frac{u^2}{e^2}.r^2}{2(n+1)} + \frac{\frac{u^4}{e^4}.r^4}{2.4(n+1)(n+3)} + \&c\right\} = 0 \end{array} \]
il y aura autant de lignes nodales circulaires
qu'il se trouvera de valeurs de $\underline{r}$ convenables
à cette equation. Les lignes nodales circulaires
\note{En tête, au milieu, « (145) » ; pas de nombre en haut à droite :
c'est le dos du feuillet de la vue 569. En tête de la première ligne,
« Mais » est barré et « Quelque » barré d'une croix, « quelles que »
écrit au-dessus ; « soient » est récrit sur « soit » ; une petite
lettre précède « Mais » dans la marge. À la fin des deux premières
séries, un signe d'égalité traversé d'un trait vertical (lu
\uncertain{$\neq$}) est suivi d'une lettre empâtée qui ressemble à un
$\alpha$ ouvert, peut-être un zéro surchargé : ni l'un ni l'autre
n'est sûr, et la phrase qui suit (« seront des quantités constantes »)
ne tranche pas. « donc » est écrit au-dessus de « ainsi », barré.
« Mais dans cette supposition » est écrit dans la marge et
l'interligne au-dessus de « Et alors », barré. Le pluriel
d'« equations » est annulé par une petite croix sur l's final : on
donne le singulier. La phrase se poursuit à la page suivante.}

\page{571}

seront de simples lignes de repos; car la supposition
d'une ligne \struck{d} circulaire d'appui entraine la condition
$A' = 0$ qui ne peut convenir qu'au seul cas où les points
du contour sont eux mêmes appuyés.

Dans le cas present des points de contour libres, il
peut donc y avoir sur la plaque circulaire deux genres
de lignes nodales, savoir \struck{des lignes d'appui qui seront}
les lignes nodales diamétrales\marginal{+ qui seront des lignes d'appui,} et des \struck{simples} \add{lignes nodales circulaires qui seront de simples} lignes de repos
qui ne satisferont pas aux conditions des extremités \struck{qui
seront les lignes nodales circulaires.}

Ces deux genres de lignes nodales peuvent exister en-
semble ou séparement sur la plaque dont le contour
est libre; ainsi les suppositions $A = 1 + \frac{\frac{u^2}{e^2}r^2}{2(n+1)} + \frac{\frac{u^4}{e^4}r^4}{2.4(n+1)(n+3)} + \&c$
et $A' = 1 - \frac{\frac{u^2}{e^2}r^2}{2(n+1)} + \frac{\frac{u^4}{e^4}.r^4}{2.4(n+1)(n+3)} - \&c$ appartiennent
au cas où il n'y a pas de lignes nodales circulaires;
\struck{et} la supposition $\beta = 0$ à celui où il n'y a pas
delignes nodales diametrales, et dans toute autre hypothese
il y aura àlafois sur la plaque circulaire
vibrante, des lignes nodales diametrales et des
lignes nodales circulaires.
\note{En tête, au milieu, « (146) », barré ; en haut à droite, à
l'encre, « 281 ». La première ligne achève la phrase de la vue 570
(« Les lignes nodales circulaires / seront de simples lignes de
repos »). La note de la marge gauche, « + qui seront des lignes
d'appui, », est appelée par la croix écrite après « diamétrales ».
L'ajout de l'interligne, « lignes nodales circulaires qui seront de
simples », finit dans la marge droite (« simples »). Dans la ligne
« qui ne satisferont pas … », le premier « qui » est traversé d'un
trait qui peut être une rature ou le trait de liaison ; on le garde, la
phrase le demande ; le « qui » final et la ligne suivante sont barrés.
« ainsi les » est récrit sur « Ainsi Les ». Le signe devant le dernier
« \&c » de $A'$ est une tache, lue « $-$ » par la série. Dans ces deux
valeurs de $A$ et $A'$ la variable est $r$ sur la page, là où la vue 570
écrivait $a$ ; et la série aux signes alternés est ici celle de $A'$, celle de la vue 570 étant celle de $A$ : on laisse le feuillet. Le bas du feuillet est blanc.}

\page{572}

Le troisième couple d'équations conviendroit au cas où les points
du contour seroient fixés, et on pourroit y satisfaire par la
condition
\[ \begin{array}{l} A\left\{1 - \frac{\frac{u^2}{e^2}r^2}{2(n+1)} + \frac{\frac{u^4}{e^4}r^4}{2.4(n+1)(n+3)} - \&c\right\} \\ \qquad + A'\left\{1 + \frac{\frac{u^2}{e^2}r^2}{2(n+1)} + \frac{\frac{u^4}{e^4}r^4}{2.4(n+1)(n+3)} + \&c\right\} = 0 \end{array} \]
Enfin le quatrième couple donne les deux conditions
\[ \begin{array}{l} A\left\{1 - \frac{\frac{u^2}{e^2}r^2}{2(n+1)} + \frac{\frac{u^4}{e^4}r^4}{2.4(n+1)(n+3)} - \&c\right\} \\ \qquad - A'\left\{1 + \frac{\frac{u^2}{e^2}r^2}{2(n+1)} + \frac{\frac{u^4}{e^4}r^4}{2.4(n+1)(n+3)} + \&c\right\} = 0 \end{array} \]
\[ \begin{array}{l} A\left\{1 - \frac{\frac{u^2}{e^2}r^2}{2(n+1)} + \frac{\frac{u^4}{e^4}r^4}{2.4(n+1)(n+3)} - \&c\right\} \\ \qquad + A'\left\{1 + \frac{\frac{u^2}{e^2}r^2}{2(n+1)} + \frac{\frac{u^4}{e^4}r^4}{2.4(n+1)(n+3)} + \&c\right\} = 0 \end{array} \]
ou celle-ci $\text{cos}(\beta\varphi + B) = 0$. Les premières ont déjà
été examinées à l'occasion du premier couple d'équations.
Quant àla seconde qui est entièrement indépendante
des valeurs de $\underline{r}$, elle montre que lorsque la surface
est périodique par rapport aux valeurs de l'angle
$\varphi$, chaque secteur formé par deux lignes d'appui
consécutives, est partagé endeux autres secteurs égaux
entr'eux par une ligne de limite analytique;
desorte que l'on peut concevoir ces secteurs
comme réellement séparés \struck{et} en deux autres
\struck{pareilles figu} surfaces de pareille figure
qui auroient chacune un de leur coté appuyé,
l'autre libre en vertu des conditions $\frac{dz}{d\varphi} = 0$ et
$\frac{d^3z}{d\varphi^3} = 0$ et \struck{les points de l'axe} dont les sons seroient
les mêmes que ceux dela plaque entière.
\note{En tête, au milieu, « (147) », non barré ; pas de nombre en
haut à droite : c'est le dos du feuillet de la vue 571. « Les
premières » est récrit sur un autre mot. Le mot barré avant « en deux
autres » est court, lu « et ». Le bas du feuillet est blanc.}

\page{573}

N.o 23. La difficulté de déterminer les valeurs de $\underline{u}$ m'empêche
d'entreprendre d'appliqu\struck{ation} \struck{à l'experience de} l'intégrale donnée
dans le N.o précédent \add{à l'experience}; mais sans s'occuper des sons qui accompagnent
chaque figure nodale, on peut toujours \struck{que} remarquer que
les fig. 104 et 109 de M\textsuperscript{r} Chladni appartiennent au cas où $\beta = 0$
\[ \begin{array}{ll}
105,\ 110^a \text{ et } 116^a \ldots\ldots & \beta = 1 \\
99,\ 106, \text{ et } 111 \ldots\ldots & \beta = 2 \\
100, \text{ et } 107 \ldots\ldots & \beta = 3 \\
101\,a \text{ et } 108 \ldots\ldots & \beta = 4 \\
\text{enfin la fig. } 102\,a \ldots\ldots & \beta = 5
\end{array} \]
Si dans l'équation differentielle
\[ \frac{u^4}{e^4}z = \ldots \left\{\begin{array}{l} \frac{1}{r^3}.\frac{dz}{dr} - \frac{1}{r^2}.\frac{ddz}{dr^2} + \frac{2}{r}\frac{d^3z}{dr^3} + \frac{4}{r^4}.\frac{ddz}{d\varphi^2} \\ - \frac{2}{r^3}.\frac{d^3z}{d\varphi^2dr} + \frac{2}{r^2}\frac{d^4z}{d\varphi^2dr^2} + \frac{d^4z}{dr^4} + \frac{1}{r^4}\frac{d^4z}{d\varphi^4} \end{array}\right. \]
on fait $r = a$, on devra trouver l'équation \add{differentielle} de la circonférence
dela plaque circulaire, ou, cequi est lamême chose, l'équation
differentielle de l'anneau élastique. Il est clair que la supposition
de $\underline{a}$ constant \struck{fait} \struck{fera} disparoître toutes les differentielles
partielles relatives à $\underline{r}$, et aussi le terme $\frac{4}{r^4}.\frac{ddz}{d\varphi^2}$ \struck{qui} dans
lequel des differentielles du même genre n'ont \struck{disparues} \add{été effacées} qu'à
l'aide des substitutions (V. N.\uncertain{o} 2). ainsi l'équation de
l'anneau se reduira à $\frac{u^4}{e^4}.z = \ldots \frac{1}{r^4}\frac{d^4z}{d\varphi^4}$
dont l'intégrale sera $z = \ldots \mathcal{A}\,\text{sin}(\beta\varphi + B)$. $\mathcal{A}$ étant une
constante arbitraire \struck{qui} rien n'empêche de \struck{faire} prendre $\mathcal{A}$
\[ \begin{array}{l} = a^\beta\left\{A\left\{1 - \frac{\frac{u^2}{e^2}a^2}{2(n+1)} + \frac{\frac{u^4}{e^4}a^4}{2.4(n+1)(n+3)} - \&c\right\}\right. \\ \qquad \left. + A'\left\{1 + \frac{\frac{u^2}{e^2}a^2}{2(n+1)} + \frac{\frac{u^4}{e^4}a^4}{2.4(n+1)(n+3)} + \&c\right\}\right\} \end{array} \]
\note{En tête, au milieu, « (148) », barré ; en haut à droite, à
l'encre, « 282 ». « N.o 23 » ouvre la page, dans la ligne. La fin de
« d'appliquation » est barrée, et la phrase ne récrit pas la fin du
verbe ; « à l'experience », barré sur la ligne, est récrit dans
l'interligne au-dessus de « précédent ». Les numéros de figures sont
ceux du traité de Chladni ; les « a » sont écrits petits, en exposant
ou sur la ligne, comme sur la page. La table est rangée en colonnes
comme sur la page, les tirets et les points de conduite rendus par
\ldots. « N.\uncertain{o} 2 » : la lettre qui suit « N. » est empâtée ; le
N.o 2 est cité ainsi à la vue 566. $\mathcal{A}$ : une capitale
cursive bouclée, distincte des $A$ et $A'$ de la même ligne. Dans la
dernière formule, le « 4 » du deuxième dénominateur est couvert d'une
tache. Le bas du feuillet est blanc.}

\page{574}

et l'équation de l'anneau se présentera alors sous laforme d'un
cas particulier de l'équation dela plaque circulaire.

Il est aisé de voir que les équations données par Euler
àlafin de son mémoire sur la lame élastique ne different
de celles que je viens de trouver, qu'en ceque les arcs
remplacent dans mes formules les ordonnées auxquelles elles
sont egales \struck{et} \add{mais} \struck{qu'} \add{que} \struck{dans celles de} cet Illustre auteur a
continué de désigner delamême manière que s'il s'agissoit
encore d'une simple lame droite.

On voit donc que soit que l'on applique immédiatement
à l'anneau l'analyse dela \struck{simple} lame \add{droite} dont les deux
extrémités sont appuyés, en se contentant, \struck{comme} \add{ainsi que} le fait
Euler, de considerer ces deux extrémités comme réunis en un
seul point. soit que l'on cherche la même equation
en employant les trois coordonnés qui doivent déterminer
la courbure, on est toujours conduit àla même équation.

Cette équation montre que les sons qui accompagnent
chaque figure nodale, doivent être proportionels \uncertain{aux}
quarrés des nombres \struck{1, 2, 3, 4 \&c.} 2, 4, 6 \&c. Euler indique
la suite des nombres 4, 9, 16, 25 \&c. mais ce grand
géomètre ne paroit pas avoir remarqué qu'il \struck{y a toujours}
ne peut y avoir qu'un nombre pair de nœuds sur
l'anneau élastique vibrant à la manière d'une lame
\note{En tête, au milieu, « (149) », non barré ; pas de nombre en
haut à droite : c'est le dos du feuillet de la vue 573. « Euler » est
écrit avec l'E en forme de L et une fin cursive (« Lober »), trois
fois. Après « élastique » un petit signe en forme de croix annule
semble-t-il un « s » final. Dans la ligne « sont egales … », « mais »
et « que » sont écrits au-dessus de « et » et de « qu' », barrés ;
« dans celles de » est barré sur la ligne. Le mot de la fin de ligne
après « proportionels » est surchargé (« à la » repris), lu
\uncertain{aux} ; « quarrés » est récrit sur un autre mot. « la
suite » : les deux mots sont récrits sur des capitales. « nœuds » est
écrit « nauds », l'œ ouvert, comme ailleurs dans ce volume. La phrase
se poursuit à la page suivante ; le bas du feuillet est blanc.}

\page{575}

Le mauvais état de ma santé m'a empêché d'étudier
le mémoire que Monsieur fourier a eu la complaisance
de me prêter, je veux pourtant lui communiquer une observation
qui me semble curieuse.

Dans le N.o d'aout, le dernier qui ait paru, M\textsuperscript{r} Poisson repond
a M\textsuperscript{r} Navier p \uncertain{438} que l'hypothese n'est point admissible
il semble ne connoître cette hypothese que j'ai donné
en 1811 que par la citation de M\textsuperscript{r} Navier.

le plus curieux est l'objection « pour s'assurer del'inexactitude
\&c. puis une equation, (que par sa forme, je crois être
applicable au point de contour mais peu importe)
\add{il dira qu'il a entendu tel ou tel restriction}
Suivant M\textsuperscript{r} Poisson l'hypothese ne seroit pas applicable
a la plaque circulaire dont les points du contour
sont assujettis a des forces normales on devoit donc
s'attendre a trouver dans son mémoire au § VIII
p. 189 ou il traite dela plaque circulaire une équation
qu'on ne pourroit déduire del'hypothese \struck{absurde}
inadmissible: je savois qu'il ne \uncertain{traiteroit} pas
d'une equation generale qui fut contraire a
la chose dont personne aumonde n'est plus convaincu
que lui mais pourtant j'ai voulu verifier.
et voici cequi je trouve.

p. 190 $\frac{d^2z}{dx^2} + \frac{d^2z}{dy^2} = \varphi$ \quad p. 191 $\frac{d^2\varphi}{dx^2} + \frac{d\varphi}{dy^2} = \frac{d^2\varphi}{dr^2} + \frac{1}{r}\frac{d\varphi}{dr}$
\note{Petit feuillet d'un autre papier, monté dans le bas d'une grande
feuille blanche ; en haut à droite, à l'encre, « 283 ». Aucune
pagination de sa main. L'écriture est plus rapide et plus pesante
que celle des pages paginées qui l'entourent ; la note est à la
première personne (« cette hypothese que j'ai donné en 1811 »), et
parle de Fourier à la troisième (« lui communiquer »). La page
(149) de la vue 574 s'arrête au milieu d'une phrase, dont la suite
n'est pas ici. « N.o d'aout » : le numéro d'août d'un périodique qui
n'est pas nommé. Le chiffre de page lu \uncertain{438} s'écrit « b 38 »,
le premier chiffre formé comme un « b », à la manière de son 4. Un
guillemet ouvrant précède « pour s'assurer » et n'est pas fermé ; la
citation s'interrompt sur « \&c. ». La ligne « il dira qu'il a entendu
tel ou tel restriction » est écrite plus petit, serrée dans
l'interligne, sous la parenthèse ; on la place où elle est. « absurde »
est barré à la fin de la ligne, et « inadmissible » suit à la ligne.
Le mot lu \uncertain{traiteroit} commence par un trait qui peut être un
« s ». Dans la seconde équation, $\frac{d\varphi}{dy^2}$ est écrit sans
exposant au $d$. La note s'arrête sur ces deux équations, au bas du
feuillet. Que le périodique soit les \emph{Annales de chimie et de
physique} et la réponse de Poisson à Navier celle de leur controverse
de 1828 est une inférence du transcripteur, que rien sur le feuillet
n'établit.}

\page{577}

droite appuyée a ses extrémités, et \struck{que chaque nombre} qu'un
nombre pair de nœuds dans l'anneau équivaut à un
nombre impair de nœuds dans la lame droite.

En traitant des conditions auxquelles sont assujettis les
points de contour des plaques libres, (V. N.o 22) j'ai eu occasion
de remarquer que\struck{\,il} dans le cas où \struck{il} n'y avoit
sur ces plaques que des lignes nodales diamétrales,
\struck{il étoit permis de} la valeur de $\underline{u}$ restoit indeterminée,
rien n'empêche donc de faire $\frac{u}{e} = \beta$, et cette valeur \struck{convient
en} conviendra en effet toutes les fois que l'on fera
abstraction dela variabilité de $\underline{r}$; les sons seront alors
proportionels aux quarrés des differentes valeurs de $\underline{\beta}$; \struck{et}
$\beta$ indiqu\struck{ant}\add{e}, comme je l'ai deja dit, le nombre des
lignes diamétrales \add{dela plaque} \struck{circulaire}, et \struck{$2\beta$} parconséquent $2\beta$ exprim\struck{ant}\add{e}
celui des nœuds de l'anneau élastique correspondant.

N.o 24. On peut voir, p. 118 du traité d'acoustique, que M\textsuperscript{r}
Chladni observe que la théorie de l'anneau donnée par
Euler n'est pas confirmée par l'experience. En effet cet
habile physicien a trouvé que les sons simples rendus
par l'anneau partagé en 4, 6, 8, \&c parties vibrantes,
sont proportionels aux quarrés des nombres 3, 5, 7 \&c.
tandis que la théorie veut que les mêmes sons soient
proportionnels aux quarrés des nombres 4, 6, 8 \&c.
\note{En tête, au milieu, « (150) », barré ; en haut à droite, à
l'encre, « 284 ». La première ligne achève la phrase laissée au bas
de la page (149), vue 574 (« … d'un nombre pair de nœuds sur
l'anneau élastique vibrant à la manière d'une lame / droite appuyée a
ses extrémités ») : le feuillet de la vue 575 s'intercale entre les
deux pages. « (V. N.o 22) » est écrit au-dessus de la ligne, après
« libres ». Dans « que il », le « il » est barré ; dans « il n'y
avoit », « il » est barré aussi. « indique » et « exprime » : la
finale « ant » est barrée dans les deux mots ; on donne ici le « e »
qui reste comme une addition, la lecture étant celle du sens.
« dela plaque » est écrit au-dessus de « circulaire », barré. Le
premier $2\beta$ de la ligne est barré. « nœuds » est écrit
« nauds ». « Chladni » s'écrit « Chladni » avec une capitale bouclée ;
« Euler », avec l'E en forme de L. Le bas du feuillet est blanc.}

\page{578}

cependant on ne peut guere douter que le cas soumis à l'experience
\struck{ne} soit veritablement celui qui repond au cas des extremités
appuyés dans la lame droite; car un caractère particulier
de cette manière de vibrer, est l'égalité parfaite des parties
séparées par des points ou lignes d'appui, et on
peut voir N.o 89 du traité d'acoustique, que l'auteur fait
une remarque expresse de cette égalité.

Pour concilier des expériences sur l'exactitude desquelles
on ne peut élever aucun doute, avec une théorie donnée
par un des géometres dont l'autorité a le plus de poids,
il faut avoir recours à la considération de quelques
circonstances \struck{\uncertain{negl}} qui n'entrent pas dans le calcul.

On peut voir section VI P. 113 du traité d'acoustique,
\struck{qu'a mesure que l'auteur aug} que plus la courbure
d'une bande droite de fer \struck{étoit augmentée} representée
fig. 37, étoit augmentée, plus aussi les sons \add{baissoient par} \struck{compar\ill{}} à
cequ'ils eussent été dans lamême manière de vibrer
dela bande droite; \struck{étoient trop baissoient.} et qu'en même
\uncertain{tems} les nœuds de vibrations situés dans la partie
infléchie, se rapprochoient de plus en plus.

Dans la fig. \uncertain{ee} qui est celle àlaquelle se rapportent
les experiences fig. 38, 39 \ldots 42, la fourche en donnant
les mêmes sons qu'une lame droite dont les extrémités
\note{En tête, au milieu, « (151) », non barré ; pas de nombre en haut
à droite : c'est le dos du feuillet de la vue 577. Le premier mot de
la deuxième ligne, barré, est récrit et lu « ne ». Le mot barré avant
« qui n'entrent » est court, lu \uncertain{negl}. Au-dessus de
« compar- », dont la fin est noyée dans une tache, est écrit
« baissoient par » ; le « à » qui suit reste. « nœuds » est écrit
« nauds ». Le mot lu \uncertain{tems}, en tête de ligne, est court et
cursif. Après « la fig. », deux petites lettres rondes, lues
\uncertain{ee}, tiennent la place d'un numéro de figure : peut-être
un blanc laissé à remplir, peut-être un chiffre mal formé ; la
référence à Chladni ne se vérifie pas sur le feuillet. Le 8 de « 38 »
est récrit. La phrase se poursuit à la page suivante ; le bas du
feuillet est blanc.}

\page{579}

eussent été appuyés, présentoit toujours deux nœuds fort voisins
et placés dans la portion courbe dela fourche; \struck{desorte que
l'un et les autres d'autant plus rapprochés les uns des autres
qu'ils etoient dans les parties} et l'intervalle entre deux nœuds
n'étoit constant que dans les parties entièrement droites des
branches de la même fourche.

Il est clair que la circonstance de la courbure considérée analytique-
ment, ne peut rendre compte des faits observés par M\textsuperscript{r} Chladni;
\struck{mais on doit remarquer} \struck{que} \add{car} dans \struck{son} le calcul, on suppose
toujours la rigidité naturelle constante dans tou\struck{tes} les points
parties du corps sonore; et la courbure considerée en elle même
ne pouvant altérer ce coëfficient, les formules doivent rester les
mêmes, soit que le corps vibrant \struck{soit} ait une forme rectiligne,
soit que le même corps affecte toute autre forme. Et en
effet lorsqu'il s'agit des instrumens à vent, l'experience
montre que la courbure n'influe pas sur les sons.
Il en est tout autrement dans les expériences de M\textsuperscript{r} Chladni
et on ne peut supposer qu'une bande de fer ait conservé
sa rigidité naturelle, lorsqu'elle s'est trouvée ployée aux
différentes formes que ce physicien lui a fait prendre.
Le résultat des expériences que je me propose d'expliquer,
montr\struck{ent} que l'effet dela torsion a été de diminuer
l'élasticité; \struck{et} les deux circonstances observées, savoir l'abaissement
des sons et le rapprochement des nœuds, concourent \add{egalement} a prouver cette diminution.
\note{En tête, au milieu, « (152) », barré ; en haut à droite, à
l'encre, « 285 ». Le feuillet est monté plus bas sur sa feuille de
support, dont le haut reste blanc. La première ligne achève la phrase
de la vue 578 (« une lame droite dont les extrémités / eussent été
appuyés »). « nœuds » est écrit « nauds ». Dans « toutes les
points », la finale « tes » est surchargée, sans que « points » soit
barré ; on garde les deux lignes telles qu'elles se lisent, « points »
en fin de ligne et « parties » en tête de la suivante. « montrent » :
la finale est barrée. « egalement » est écrit au-dessus de « a
prouver » ; « cette diminution » est serré au bord droit du feuillet.
Le texte s'arrête à la dernière ligne du feuillet, sur un point
incertain.}

\page{580}

En effet, toutes choses égales d'ailleurs, le son est d'autant plus aigu
que l'élasticité est plus grande; ainsi lorsque tout est constant,
excepté ce coëfficient, le son ne peut baisser sans qu'\struck{il} \add{l'élasticité} diminue.
Si le son reste le même, quoique l'élasticité diminue, \struck{tout}
lorsque le poids et l'épaisseur seront constans, il faudra
que l'étendue dela partie vibrante soit aussi diminuée.
On peut conclure delà que si le changement d'élasticité
n'a lieu que dans quelques points d'une lame élastique, il
doit arriver àlafois deux choses, l'une que le son baisse,
l'autre que les parties vibrantes séparées par des nœuds
de vibrations, soient d'autant moins grandes que leur élasticité
est plus diminuée, c'est adire qu'elles se trouvent dans des
portions dela fourche plus fortement infléchies. et
\add{rien n'est plus conforme en effet à} \struck{On voit donc que c'est précisément là} ceque M\textsuperscript{r}
Chladni a observé dans les differens cas du mouvement
des fourches.

Les mêmes principes peuvent servir à prévoir ceqi doit
arriver lorsque la torsion, aulieu de se faire dans un espace
circonscrit pris sur la bande élastique, a lieu uniformément
dans toutes les parties dela même bande. Alors, il ne peut
se manifester qu'un seul des deux effets qui se montrent
dans le cas des fourches; car l'élasticité qui est proportionnelle
àla courbure, est également constante dans l'étendue de la
bande lorsque sa courbure est nulle, comme dans la bande
\note{En tête, au milieu, « (153) », non barré ; pas de nombre en
haut à droite : c'est le dos du feuillet de la vue 579. « l'élasticité »
est écrit au-dessus de « qu'il », dont le « il » est barré. Le pluriel
de « fourches », dans « portions dela fourche », est annulé par une
petite croix sur l's final : on donne le singulier. La ligne « rien
n'est plus conforme en effet à » est écrite dans l'interligne, plus
petit, au-dessus de la ligne barrée qu'elle remplace. « nœuds » est
écrit « nauds ». « ceqi doit » : ainsi sur la page. La phrase se
poursuit au-delà de ce lot ; le bas du feuillet est blanc.}

\end{document}
